% SOURCE_AUTHORITY: sga5_fr_workpass.tex, lines 9868--9901 (LF-normalized unit).
% SOURCE_SHA256: 791F4EFFC5E02832D5D77ED03518C8156D6F07E4C8238B03545DB93D883FBB28
Aplicando $g_*$ a ambos os membros de 4.1.3, obtém-se, portanto,
\[
g_*(u)=a_{d-1}.
\tag{4.1.6}
\]
Por outro lado, aplicando $f_*$ a ambos os membros de 4.1.1, obtém-se, tendo em conta 4.2,
\[
i_*(1_Y)=i_*g_*(u).
\tag{4.1.7}
\]
Mas o homomorfismo $i_*$ em questão é injetivo, de onde
\[
a_{d-1}=g_*(u)=1_Y,
\]
que termina com demonstração.

\begin{statement}{Observaciones 4.4}
(a) O ingrediente essencial do que precede é o teorema da pureza. Uma vez disponível em toda generalidade, a frase 4.1 será válida com $Y$ regular de código puro $d$ em $X$, e $X$ excelente e regular em um ambiente de $Y$. Nessa forma, o resultado é estabelecido em qualquer caso em característica $0$, graças a (SGA 4 XIX).

b) Quando $S$ é o espectro de um corpo separadamente fechado e $X$ e $Y$ são quase projetivos e suaves sobre $S$, MUMFORD estabeleceu a afirmação analógica no anel de Chow de $Y$ (cf. 9.2).
\end{statement}

\begin{statement}{Corolário 4.5}
A partir das hipóteses de 4.1, sejam $x$ e $y\in A^\bullet(Y)$. Tem
\[
i_*(x)i_*(y)=i_*(xy\,c_d(\check N)).
\]
\end{statement}

\emph{Demonstração.} A fórmula de projeção para o morfismo $i$ mostra que
\[
i_*(x)i_*(y)=i_*(x\,i^*i_*(y)),
\]
de onde segue-se a conclusão graças a 4.1.
